\section{Conclusions, baseline, and research scope}\label{sec:scope}

The strongest new mathematical step in this continuation is an exact
finite-threshold recurrence for one long twist region. The existing
support-counting formula says that the chain dimension grows affinely in a
selected exponent. We pass from that statement about the number of
generators to a statement about the differential and hence the actual
homology. This distinction matters: a repeated vector space can carry
different differentials, and chain-size growth alone does not determine
homology growth.

The implementation work has two complementary purposes. The recurrence
avoids constructing an arbitrarily long repeated part. The streamed backend
reduces retained storage when the finite complex still needs to be evaluated.
Direct braid-source structural bounds avoid expanding a binary exponent
merely to rediscover a certificate already available from elementary counts.
Each change has an explicit interface and an independently stated scope.

\subsection{What is already in ProveIt}

The baseline is the repository's main revision
\code{ea2abcb115aaa58f0b193ce1e045c2def983e1e6}, retrieved on 8 October 2026.
The immediately relevant integrated work comprises report 10 on structural
and algebraic improvements, report 11 on twist compression, report 12 on
component algebra, and their synthesis updates~\cite{repo}.
The complete \code{fast/} Python source, tests, examples, and build metadata
at this revision accompany the package as a reference snapshot.

The current recognizer already has much more than a naive crossing cube.
Its pipeline contains checked diagram construction, braid certificates
including the at-most-three-strand endpoint, signed Seifert-graph criteria,
diagram simplification, visible connected-sum and interlacement
factorization, modular Alexander and Jones obstructions, and exact
Khovanov scanning. Optional paths provide a twist macro backend and
component-algebra contractions. Existing scan machinery includes local
gluing, unit cancellation, bit-packed linear algebra, and choices of scan
order. Those features are baseline capabilities, not contributions of this
continuation.

In particular, report 11 explicitly proposed degree-streamed rank
computation and required the state generator and map caches to be included
in its memory analysis. We implement that proposal rather than presenting
the rank identity underlying it as a new mathematical discovery.
Report 11 also established, for supplied signed runs with lengths \(m_j\),
the support polynomial
\begin{equation}\label{eq:support}
\N=\sum_{S\subseteq[t]}2^{c(S)-1}\prod_{j\in S}m_j
\end{equation}
and a restricted-family complexity bound
\begin{equation}\label{eq:prior-bound}
\poly(n)\,2^{O(t\log(2+n/t))},\qquad n=\sum_jm_j.
\end{equation}
Here \(c(S)\) is the number of circles after choosing the turnback exactly
in the boxes in \(S\). The run count \(t\) is the count in the supplied
presentation. It is not a proven small invariant of every knot.

\begin{center}
\small
\begin{tabular}{p{0.43\linewidth}p{0.48\linewidth}}
\toprule
Result or capability & Status in this continuation\\
\midrule
Two-strand twist normal form; Khovanov unknot detector &
Established literature, used with explicit attribution.\\
Support polynomial and \(t=O(\log n)\) bound &
Already integrated in the pinned ProveIt baseline.\\
One-reference signed homology recurrence and slope extraction &
Proved here as a continuation of that model; implemented.\\
Slope as reduced homology of the turnback closure &
Proved using established dot-sliding and local contraction machinery.\\
Degree-streamed ranks, bounded caches, early lower bound &
Implemented with exact accounting and regression checks.\\
Structural certificates on binary run input &
An encoding improvement to existing sufficient conditions.\\
New signature packing decision class &
None: the proposed bound is proved dominated.\\
Unrestricted quasi-polynomial recognition &
Not established by this work.\\
\bottomrule
\end{tabular}
\end{center}

\subsection{Current literature and the asymptotic target}

The mathematical local input is Bar--Natan's tangle category and
cancellation calculus~\cite{bn05,bn07}, with integer-twist complexes
described explicitly by Thompson~\cite{thompson}. The rank-one decision
criterion is due to Kronheimer and Mrowka~\cite{km}.
Twist stabilization and compact tangle descriptions have substantial
precedent, including the immersed-curve framework of Kotelskiy, Watson,
and Zibrowius~\cite{kwz}. The finite threshold and one-reference extraction
below should be assessed as a precise, executable continuation of the
repository's model. We make no general priority claim for stability under
twisting.

Kelom\"aki and Sch\"utz's 2026 preprint supplies an especially useful
boundary for complexity arguments: their Theorem 1.2 gives exponential
scanning behavior on certain three-strand braid families, whereas their
Theorem 1.3 gives a specialized polynomial-time computation for all
three-strand braids~\cite{ks26}. A bounded number of boundary matching
types therefore cannot, by itself, bound the number of algebraic summands.
The specialized theorem also prevents a misleading benchmark claim:
three-strand examples are useful algebraic checks, but do not demonstrate
new difficult instances for the integrated recognizer.

The hierarchy route must be described separately. Lackenby's July 2026
preprint develops incompressible surfaces, boundary patterns, and unknot
recognition. Section 9 bounds iterations while explicitly leaving costs
of steps insufficiently specified for a complete runtime estimate;
Proposition 9.1 gives \(L(g+1)^L\) iterations in its hierarchy and pattern
parameters~\cite{lackenby26}. Other results provide compact surface
certificates and algorithms polynomial in handle data and logarithmic
normal weight. Those facts support a serious alternative research
programme, but they do not supply the missing cost theorem for the present
Khovanov implementation. Certificate existence, certificate verification,
certificate search, and repeated compressed operations remain different
tasks.

We also inspected the user-suggested \code{openai/math} catalogue and
repository guidance~\cite{openaimath}. Its manuscripts have varying
verification status. No theorem from that collection is needed for the
proofs here; a superficially similar quasi-polynomial result in a different
subject would not remove the particular representation and rank barriers
identified below. The useful methodological lesson is to keep dependencies,
proof status, and machine-checkable artifacts explicit.

\subsection{Four sizes that must be kept distinct}

Let
\begin{equation}\label{eq:input}
\beta=\sigma_{i_1}^{e_1}\cdots\sigma_{i_t}^{e_t},\qquad
1\le i_j<b,\quad e_j\in\Z\setminus\{0\}.
\end{equation}
Put \(m_j=|e_j|\), \(n=\sum_jm_j\), and define the encoded input size,
up to a fixed choice of delimiters, by
\begin{equation}\label{eq:encoded}
\lambda=1+\lceil\log_2(b+1)\rceil+
\sum_{j=1}^t\bigl(1+\lceil\log_2(b+1)\rceil+
                         \lceil\log_2(m_j+1)\rceil\bigr).
\end{equation}
An expanded braid word has length \(n\). A run-encoded braid can have
\(n\) exponential in \(\lambda\). A planar diagram normally presents its
crossings explicitly and therefore has a different input model.

There are also the number \(S=\prod_j(m_j+1)\) of macro states, the total
reduced chain dimension \(\N\), and the bit length of the requested output.
A scalar formula for \(\N\) does not compute a differential rank. A
compact formula for a homological profile is a different output from an
explicit dictionary containing each degree. We charge for the chosen
representation throughout.

The closure of a braid is a knot exactly when its strand permutation is a
single cycle. A run contributes the adjacent transposition only when its
exponent is odd. Thus original component count is computable without
expanding any run. For a one-component \(b\)-braid every generator index
\(1,\ldots,b-1\) occurs; omitting index \(i\) separates the first \(i\)
strands from the others. This simple fact gives \(b-1\le t\) and will also
control the reference size in the dominant-run theorem.

\subsection{The characteristic-two exact endpoint}

For a knot \(K\), the endpoint used here is
\begin{equation}\label{eq:detector}
K\text{ is the unknot}\quad\Longleftrightarrow\quad
\dim_{\F}\rKh(K;\F)=1.
\end{equation}
The implication from rank one follows from the rational rank-one detector
of Kronheimer--Mrowka and the universal coefficient theorem: characteristic
two dimension is at least the rational dimension, and reduced homology is
nonzero. The unknot has one reduced generator over either field.
All calculations retain a fixed reduction point. For links we compute
reduced homology at that point, but do not apply a knot-only decision
criterion to them.

Resource exhaustion has a precise meaning. A mathematical exact algorithm
without fixed caps terminates on finite inputs. Its resource-bounded
implementation can return \code{UNKNOWN}. Such a result is neither
evidence of knottedness nor evidence of triviality. The same distinction
applies when a computation has established only a lower bound on the
homology dimension.

\section{The local complex and its global counts}\label{sec:local}

We use \(A=\F[X]/(X^2)\), with
\[
\Delta(1)=1\otimes X+X\otimes1,\qquad \Delta(X)=X\otimes X.
\]
A resolution with \(c\) circles has reduced space of dimension
\(2^{c-1}\): the circle containing the fixed mark carries \(X\), and each
other circle carries \(1\) or \(X\).
Let \(I\) denote the identity two-strand matching and \(E\) the matching
joining the two top ends and joining the two bottom ends. Names such as
cap and cup depend on drawing orientation; the chosen arcs are always
identified as the top and bottom arcs of \(E\).

\begin{proposition}[Established local normal form]\label{prop:normal}
After forgetting quantum grading, a positive run of length \(m\) is
represented by
\[
I\xrightarrow{s}E\xrightarrow{L}E\xrightarrow{L}\cdots
\xrightarrow{L}E,\qquad L=X_{\rm top}+X_{\rm bottom},
\]
in degrees \(0,1,\ldots,m\). A negative run is represented by its dual
complex in degrees \(-m,\ldots,-1,0\), with \(I\) in degree zero, the final
saddle into \(I\), and the same internal dot sum.
The replacement is valid inside any fixed exterior tangle.
\end{proposition}

\begin{proof}
This is the characteristic-two, quantum-forgotten specialization of the
integer-twist normal form~\cite[Proposition 4.1]{thompson}. Tensoring with
one additional crossing gives an \(E\)-circle summand. Delooping that circle
and cancelling an invertible saddle block gives the next zig-zag term.
The local contraction formula is the categorical Schur complement of
Bar--Natan~\cite{bn07}. Over the integers the internal maps have alternating
signs; the dot sum and difference agree over \(\F\). Mirroring reverses
degrees and dualizes saddles. These operations are relative to the
boundary, which justifies the stated exterior compatibility.
\end{proof}

In the tensor product of run complexes, a state has coordinates
\(0\le k_j\le m_j\). Its degree is \(\sum_j\operatorname{sign}(e_j)k_j\).
Its geometry depends only on the support \(\{j:k_j\ne0\}\), while its
differentials also distinguish the first saddle from internal dot maps.
For a fixed support \(S\), the number of coordinate choices is
\(\prod_{j\in S}m_j\). This proves \eqref{eq:support}.
The independence of geometry from the nonzero coordinate is also the
key to the stronger homology result below.

For later accounting define
\begin{equation}\label{eq:layers}
d_h=\dim C^h,\qquad r_h=\rank(\partial_h:C^h\longrightarrow C^{h+1}),
\qquad \N=\sum_h d_h.
\end{equation}
The exact formula
\begin{equation}\label{eq:betti}
\beta_h=\dim H^h=d_h-r_{h-1}-r_h
\end{equation}
uses both incident boundary ranks. Computing only chain dimensions, only
one boundary rank, or a rank bound in the wrong direction cannot justify
an exact Betti number.

Two convenient coarse quantities are
\[
Q=\sum_h d_hd_{h+1},\qquad
T_{\rm elim}=
O\!\left(\N+\sum_h d_h d_{h+1}\min(d_h,d_{h+1})\right)
\]
for elimination once the actual columns are supplied. Generating states
and local maps is a separate cost; it is included explicitly in the
algorithm bounds below. The \(\N\) term also charges for visiting zero
columns and end degrees.
\(Q\) is a dense bit envelope, not a count of actual nonzero entries.
The displayed elimination term is a conservative bit-operation bound for
ordinary characteristic-two Gaussian elimination. It is enough for
asymptotic upper bounds; it is not a claim that Python attains an optimal
matrix multiplication exponent.
